% =====================================================================
\section{Limitations and Conclusion}\label{sec:results}
\subsection{What was measured}
\sheetfig{s00_cover_kpis}
\begin{table}[h]
\centering\small
\caption{Quantities taken from the program output of this report's build (not from memory).}
\label{tab:measure}
\begin{tabular}{@{}p{8.7cm}p{7.2cm}@{}}
\toprule
Quantity & Value (source)\\\midrule
Leaf parts in the scene, in \nObjects{} export objects & \nLeaves{} (\code{scene\_export.json}, Table~\ref{tab:parts})\\
Parts with emissive $>0$ & \nEmissive{} (\code{scene\_export.json})\\
Track links of the tank & \nTrackLinks{} (export; T80)\\
Mesh memory, indexed vs non-indexed & \memIndexed{} B vs \memNonIndexed{} B (Table~\ref{tab:memory})\\
Automated tests & \nTests{} run, \nPass{} pass, \nFail{} fail (\code{test\_results.json})\\
Terrain & 49\,921 vertices, 97\,440 triangles, 36\,m radius (T13, T61)\\
Bunker clear height / walk eye height & 2.60\,m / 1.7\,m (T62, T51)\\
Missile: first launch request after lock & 1.53\,s (T108); detonation 0.50\,m from the spot (T100)\\
Shell range on level ground & 203.8\,m, $-0.24$\,\% from theory (T101)\\
Debris: pieces at rest within 8\,s & 23 of 23 (T106)\\
180\,s headless run: UAV max radius, min clearance, rover hidden & 15.57\,m, 5.62\,m, 61\,\% of frames (T45)\\
Full base, 120\,s simulated & 9 controllables (7 possessable), finite state (T113)\\
\textbf{FPS / ms per frame on the RTX 4070} & \sNM; command in Section~\ref{sec:platform}\\\bottomrule
\end{tabular}
\end{table}

Up to v5 no frame time was measured, and the EPIC preset was exercised only on llvmpipe; the v6 and v7 runs on the RTX 4070 are covered in Section~\ref{sec:v6results}. The GPU timer labels in Section~\ref{sec:pipeline} are what \cli{benchmark} reports on the RTX machine.

\subsection{Limitations}
\begin{itemize}
  \item \textbf{Verification environment.} Up to v5, every capture and test run used Mesa llvmpipe, and the visual tuning (explosion brightness, interior exposure) was done on llvmpipe images. From v6 on, captures and tests run on the RTX 4070 through the WSL2 D3D12 driver; Section~\ref{sec:v6results} gives those results.
  \item \textbf{Windows.} The executable is cross-compiled with zig and has never been run natively.
  \item \textbf{Soft particles.} The effect shader supports a soft depth fade but the renderer passes \code{uHasSceneDepth=0}, so smoke clips hard against geometry.
  \item \textbf{Combat scope.} Hit tests use spheres (rovers, radius 0.4\,m) and boxes (static solids), not part meshes. The player's machines have no damage state other than camouflage scorching, and nothing can hit a UAV.
  \item \textbf{Exporter.} It labels surfaces 12--18 as ``Other'' and the terrain mesh function string with radius 20 (stale labels).
\end{itemize}

\subsection{Conclusion}
ThermalSense is a night military base built within the lab's rules: \nLeaves{} parts, all unit cubes, spheres, cylinders and cones with $T\,R\,S$ transforms,
share six indexed meshes. It has a walkable underground bunker with eased blast doors, a base apron with an animated flag, windsock and gate, and machines with
non-trivial motion (arc-length track links, skid steer with a stabilised turret, a cyclic rotor with a spring-damper attitude, a spooling Gatling gun, three radar modes). Two UAVs render the scene as heat, a computer-vision
pipeline and a multi-target tracker turn it into tracks, and a squad coordinates proportional-navigation strikes on armoured, infiltrating rovers (Chapters~\ref{sec:vision}, \ref{sec:squad}); the game logic is deterministic and covered by \nTests{} CPU tests, \nPass{} of which pass.
The remaining work is measurement, not construction: FPS on the RTX 4070 and a native Windows run. The audit in Appendix~\ref{app:audit} lists every item and its evidence.
